\documentclass[12pt,a4paper]{article}
\usepackage[utf8]{inputenc}
\usepackage[margin=1in]{geometry}
\usepackage{amsmath}
\usepackage{amssymb}
\usepackage{graphicx}
\usepackage{xcolor}
\usepackage{hyperref}
\usepackage{fancyhdr}
\usepackage{listings}

\pagestyle{fancy}
\lhead{Estratégia Adaptativa Multi-Fator}
\rhead{Etapa 2: Modelo}
\cfoot{\thepage}

\lstset{
    language=Python,
    basicstyle=\ttfamily\small,
    breaklines=true,
    keywordstyle=\color{blue},
    commentstyle=\color{gray},
    stringstyle=\color{red},
    backgroundcolor=\color{white}
}

\title{\textbf{ESPECIFICAÇÃO TÉCNICA DO MODELO}\\[0.5cm]
\textit{Adaptive Multi-Factor Regime-Switching Implementation}}
\author{Quantitative Research}
\date{\today}

\begin{document}

\maketitle

\section{Arquitetura Geral}

O modelo é composto por 4 módulos integrados:

\[
\boxed{\text{Data} \rightarrow \text{Factor Computation} \rightarrow \text{Regime Detection} \rightarrow \text{ML Predictions} \rightarrow \text{Optimization}}
\]

\begin{figure}[h]
\centering
\includegraphics[width=0.95\textwidth]{reports/build/model_pipeline.png}
\caption{Pipeline operacional do modelo adaptativo multi-fator}
\end{figure}

\section{Módulo 1: Computação de Fatores}

\subsection{Fator VALUE}

\[
\text{VALUE}_i = \frac{\text{Book Value}_i}{\text{Preço}_i}
\]

Assets com elevado book-to-market são undervalued e oferecem maior retorno esperado (valor inverte o ratio).

\subsection{Fator MOMENTUM}

\[
\text{MOM}_i = \sum_{j=1}^{252} r_{i,t-j}
\]

Retorno acumulado nos últimos 12 meses. Captura persistência de tendência (winner stocks continuam subindo).

\subsection{Fator QUALITY}

\[
\text{QUALITY}_i = \text{ROE}_i = \frac{\text{Net Income}_i}{\text{Shareholders' Equity}_i}
\]

Empresas com elevado ROE apresentam melhor eficiência operacional e crescimento superior.

\subsection{Fator LOW-VOLATILITY}

\[
\text{LOWVOL}_i = \frac{1}{\sigma_i + \epsilon}
\]

Inverso da volatilidade realizada. Captura prefência por ativos menos voláteis (risk-averse investors).

\subsection{Fator PROFITABILITY}

\[
\text{PROF}_i = \frac{\text{EBIT}_i}{\text{Total Assets}_i}
\]

Retorno sobre ativos operacionais. Indica qualidade intrínseca da empresa.

\subsection{Fator GROWTH}

\[
\text{GROWTH}_i = \frac{\text{Revenue}_{i,t} - \text{Revenue}_{i,t-1}}{\text{Revenue}_{i,t-1}}
\]

Taxa de crescimento de receita. Empresas crescentes possuem maior múltiplo de valuation.

\subsection{Fator MEAN-REVERSION}

\[
\text{MEANREV}_i = -\sum_{j=1}^{63} r_{i,t-j}
\]

Negativo do retorno acumulado últimos 3 meses. Captura reverter-à-média (quem cai muito tende a subir).

\subsection{Normalização de Fatores}

Todos os fatores são padronizados:

\[
\text{factor}_{\text{normalized}} = \frac{\text{factor} - \text{mean}(\text{factor})}{\text{std}(\text{factor}) + \epsilon}
\]

Garante que todos os fatores tenham peso comparável na otimização.

\section{Módulo 2: Detecção de Regime via HMM}

\subsection{Especificação do Modelo}

Hidden Markov Model com 3 estados latentes:

\[
p(s_t, s_{t-1}, \ldots, s_1 \mid r_1, \ldots, r_T) \propto p(r_t \mid s_t) \cdot p(s_t \mid s_{t-1})
\]

\subsection{Observações (Features)}

O HMM usa dois sinais:

\begin{align}
X_{1,t} &= \frac{1}{252}\sum_{j=0}^{251} r_{market,t-j} \quad \text{(média móvel de retornos)} \\
X_{2,t} &= \sqrt{\frac{1}{252}\sum_{j=0}^{251} r_{market,t-j}^2} \quad \text{(volatilidade móvel)}
\end{align}

Ambos são padronizados antes de passar para HMM.

\subsection{Transições de Estado}

Matriz de transição $P$ modela probabilidade de mudança de regime:

\[
P = \begin{pmatrix}
p_{11} & p_{12} & p_{13} \\
p_{21} & p_{22} & p_{23} \\
p_{31} & p_{32} & p_{33}
\end{pmatrix}
\]

Onde $p_{ij} = \Pr(s_t = j \mid s_{t-1} = i)$.

\subsection{Viterbi Algorithm}

Para recuperar a sequência mais provável de regimes:

\[
\hat{s}_{1:T} = \arg\max_{s_{1:T}} \log p(s_{1:T} \mid r_{1:T})
\]

Usado para identificar o regime corrente a cada dia.

\section{Módulo 3: ML Ensemble para Previsão}

\subsection{Gradient Boosting Regressor}

Treinado com hiperparâmetros otimizados:
\begin{itemize}
    \item $n\_estimators = 200$ (200 árvores)
    \item $learning\_rate = 0.05$ (taxa de aprendizado moderada)
    \item $max\_depth = 5$ (profundidade controlada contra overfitting)
    \item $subsample = 0.8$ (80\% de dados por iteração)
\end{itemize}

Captura relações não-lineares como:
\[
r_{i,t} = f(X_1, X_2, X_3) + \varepsilon_t
\]

onde $f$ é aproximada por soma de árvores de decisão.

\subsection{Ridge Regression}

Minimiza:
\[
\min_{\beta} \|y - X\beta\|_2^2 + \alpha \|\beta\|_2^2
\]

Com $\alpha = 1.0$. Oferece estabilidade numérica e evita multicolinearidade entre fatores.

\subsection{Ensemble Voting}

Predições finais combinam ambos os modelos com pesos dependentes do regime:

\[
\hat{r}_{i,t} = \begin{cases}
0.6 \cdot \hat{r}_{GB} + 0.4 \cdot \hat{r}_{Ridge} & \text{se } s_t = \text{Bull} \\
0.4 \cdot \hat{r}_{GB} + 0.6 \cdot \hat{r}_{Ridge} & \text{se } s_t = \text{Bear} \\
0.5 \cdot \hat{r}_{GB} + 0.5 \cdot \hat{r}_{Ridge} & \text{se } s_t = \text{HighVol}
\end{cases}
\]

Intuição: Em bear markets, preferimos estabilidade (Ridge). Em bull markets, capturamos não-linearidades (GB).

\section{Módulo 4: Otimização de Portfolio}

\subsection{Formulação}

Mean-Variance com restrições dinâmicas:

\[
\max_w \quad w^T \hat{\mu} - \frac{\lambda}{2} w^T \Sigma w
\]

Sujeito a:
\begin{align}
\sum_i |w_i| &\leq L(s_t) \\
|w_i| &\leq w_{\max}(s_t) \\
\|w - w_{t-1}\|_1 &\leq \delta_{\max}
\end{align}

\subsection{Leverage Dinâmico}

\[
L(s_t) = \begin{cases}
1.3 & \text{Bull (confiança alta)} \\
0.7 & \text{Bear (reduzir riscos)} \\
1.0 & \text{HighVol (posição neutra)}
\end{cases}
\]

\subsection{Peso Máximo por Ativo}

\[
w_{\max}(s_t) = \begin{cases}
0.08 & \text{Bull} \\
0.05 & \text{Bear} \\
0.06 & \text{HighVol}
\end{cases}
\]

Evita concentração excessiva. Mais apertado em bear markets.

\subsection{Reescalamento para Target Volatility}

Portfolio é ajustado para atingir volatilidade alvo de 8\% a.a.:

\[
w_{\text{final}} = w_{\text{optimized}} \cdot \frac{\sigma_{\text{target}}}{\sigma_{\text{portfolio}}}
\]

Garante que risco é controlado independentemente do retorno esperado.

\section{Pipeline Completo}

\begin{enumerate}
    \item \textbf{Input:} historical returns, prices and fundamentals.
    \item Compute factors using data up to day $t-1$.
    \item Normalize factors to zero mean and unit standard deviation.
    \item Run HMM inference to detect regime $s_t$.
    \item Predict returns using the ML ensemble.
    \item Estimate covariance matrix $\Sigma_t$.
    \item Optimize portfolio weights $w_t^*$.
    \item Execute trades to rebalance to $w_t^*$.
    \item Monitor realized returns and update metrics.
\end{enumerate}

\section{Validação Cruzada}

\subsection{Walk-Forward Validation}

Para evitar look-ahead bias:

\begin{enumerate}
    \item Treinar ML em período $[t-252, t-20]$ (1 ano de dados, deixando 20 dias de teste)
    \item Validar em período $[t-20, t]$
    \item Mover janela para frente 1 dia
    \item Repetir para todo o período de backtest
\end{enumerate}

\subsection{Stress Testing}

Modelo testado em cenários extremos:
\begin{itemize}
    \item 2008 Financial Crisis (pior caso)
    \item 2020 COVID Crash (teste de resiliência)
    \item 2000 Dot-Com Bubble (teste de bear market)
\end{itemize}

\section{Resumo de Complexidade}

\begin{center}
\begin{tabular}{|l|c|}
\hline
\textbf{Componente} & \textbf{Complexidade} \\
\hline
Computação de Fatores & $O(n \cdot p)$ (linear) \\
HMM Inference (Viterbi) & $O(T \cdot k^2)$ (quadrática em estados) \\
ML Training (GBM) & $O(n \log n)$ por iteração \\
Portfolio Optimization & $O(p^3)$ (inversão de matriz) \\
\hline
\end{tabular}
\end{center}

Tempo total de execução diária: $\sim 2-5$ segundos para 50 ativos.

\end{document}
